% =====================================================================================
% sections/s1_intro.tex -- Introduction (body only).
% The contributions are stated without numbers; the numbers are in Sections 4 to 7 and,
% once, in Section 8.1.  No figures, tables or new computation.
% =====================================================================================
\section{Introduction}\label{s1_intro:sec}

A physics-informed neural network (PINN) represents the solution of a partial
differential equation by a neural network whose weights minimise the squared residual of
the equation at collocation points, plus the misfit to the initial and boundary
data~\citep{raissi2019}; neural trial functions for differential equations go back at least
to \citet{lagaris1998}. The Deep Ritz method~\citep{eyu2018} minimises instead a
Monte-Carlo estimate of the variational energy of an elliptic problem, with the boundary
condition imposed by a penalty. Both methods need no mesh, only automatic differentiation;
they are short to implement~\citep{lu2021} and extend formally to inverse
problems~\citep{raissi2019} and to high dimension~\citep{sirignano2018}, where other
neural solvers exist as well~\citep{han2018}; see the
reviews~\citep{karniadakis2021,cuomo2022}.

Their difficulties are equally well documented. The terms of a PINN loss can be badly
unbalanced~\citep{wangteng2021}; a neural-tangent-kernel analysis shows that they are then
learnt at different rates and exhibits a two-mode wave problem on which the unweighted
method fails~\citep{wangyu2022}, which we reuse as problem \PWone. PINNs that succeed on simple
linear equations fail as the coefficients grow, for reasons of optimisation rather than
expressivity~\citep{krishnapriyan2021}, and time-dependent problems may need a loss that
respects causality~\citep{wangsankaran2024}. Quasi-Newton optimisation has been part of
the method since \citet{raissi2019}; \citet{rathore2024} measure the benefit of running
\Adam\ and then \LBFGS\ and relate it to the ill-conditioning of the loss. The placement of collocation
points has been compared systematically by \citet{wu2023}, with residual-based refinement
and importance sampling as adaptive variants~\citep{lu2021,nabian2021}; \citet{daw2023}
attribute failures of PINNs to a failed propagation of the solution from the initial and
boundary points into the interior. Against classical
solvers, PINNs did not beat the finite element method in solution time and accuracy in the
study of \citet{grossmann2024}; \citet{mcgreivy2024} show that weak baselines and selective
reporting have made machine-learning solvers for fluid-related equations look better than
they are; and much larger benchmarks than ours exist~\citep{hao2023pinnacle}. On the theoretical side, PINN
minimisers converge for linear elliptic and parabolic problems as the number of collocation
points grows~\citep{shin2020}, and a boundary penalty is known to change the problem being
solved~\citep{babuska1973,muller2022}.

This article asks what the two methods deliver on the diffusion, wave and Laplace equations
when every error is measured against an exact solution, every run is seeded, comparisons are
paired and a classical solver is run alongside, and it contributes the following.
\begin{enumerate}[label=(C\arabic*),leftmargin=*]
\item Eight test problems, each stated so that it has a unique solution in a stated class:
the unique classical solution or, for the three-dimensional problem with discontinuous
data, the unique solution that attains the datum in mean square. The exact solutions are
given with proofs (Section~\ref{s2_problems:sec}; Supplementary
Sections~\ref{S1:sec} and~\ref{sA_proofs:sec}).
% src: research_benchmark/results/benchmark_runs.csv (350 rows = 175 networks x 2 arms)
\item A seeded benchmark of 175 networks on five low-dimensional problems, with five
collocation strategies and two optimiser arms of equal length, reporting held-out relative
errors; for the heat equation the error is also resolved in time, which shows how a small
space--time error hides a large error at late times (Sections~\ref{s3_methods:sec}
and~\ref{s4_lowdim:sec}).
\item The effect of replacing the second half of the \Adam\ iterations by \LBFGS, measured
pair by pair and also against the best logged \Adam\ iterate
(Section~\ref{s4_lowdim:sec:optimiser}).
\item A paired comparison of the collocation strategies, with a control experiment that
traces the one large penalty after the \LBFGS\ phase, that of a cell-centred grid on the
heat problem, to the absence of residual points next to the initial line
(Section~\ref{s4_lowdim:sec:strategy}).
\item Longer runs and a compatible-data control on the two- and three-dimensional problems,
which show what limits them; runs of up to 41\,000 iterations showing that the two-mode wave
problem is not solved by the unit-weight PINN at the network size used, with the trained
networks projected onto the two modes; and a comparison of all five problems with
second-order finite differences in accuracy and time (Section~\ref{s5_hard:sec}).
\item A comparison of the PINN and the Deep Ritz method on the unit cube in dimension $d=2$
to $20$, in accuracy per iteration, in cost per iteration, at equal compute and at two
budgets, with the penalty bias of Deep Ritz computed exactly
(Sections~\ref{s6_highdim:sec:equal-its}--\ref{s6_highdim:sec:weights}).
\item On the unit cube in any dimension $d$, two-sided bounds of exact order $d^{-1/2}$ for
the squared norm of the harmonic extension from $L^2(\dOm)$, with surface measure, to
$L^2(\Om)$, which is the inverse of the first biharmonic Steklov eigenvalue. With them the
$L^2$ error of any network is bounded by the two terms of the PINN loss with a constant
explicit in $d$, and the penalty bias of Deep Ritz by a sharper constant. The bound is
evaluated under a pre-registered protocol on the trained networks of the dimension study
and, without pre-registration, on networks retrained with separately written code
(Section~\ref{s6_highdim:sec:bound}).
\item A pre-registered test, for both methods, at equal iterations and at equal CPU time,
with two pre-registered controls for Deep Ritz (lift3u, rep3) and an exploratory centred-input
arm for both methods, of two cheap remedies for the plateau at
$d=20$, where neither method improves on the best affine approximation of the solution: a
fixed affine part computed by one linear solve, and coordinate-wise Legendre input features
(Section~\ref{s6_highdim:sec:remedies}).
\item Two implementation pitfalls demonstrated on minimal, released examples: a PDE
residual made identically zero by differentiating with respect to a slice of the input and
replacing the missing derivative by zeros, and a
boundary-value problem with too few conditions, on which training reaches a small loss with
a different function for every seed (Section~\ref{s7_pitfalls:sec}).
\end{enumerate}

The study combines a dimension-explicit stability result with controlled measurements of
established neural PDE solvers. The two-sided bound in (C7) is the theoretical contribution;
the experiments connect optimiser choice, sampling, error reporting and computational cost
within one reproducible protocol.
% src: refs.bib (rathore2024, annote: Sec. 2.2 and Table 1 of arXiv:2402.01868v2, checked 2026-10-01)
The gain from an \LBFGS\ phase after \Adam\ was measured by \citet{rathore2024}; what is
added here is the same measurement, paired by seed, on the heat and Laplace problems at a
small budget and against an untuned \Adam\ (Section~\ref{s4_lowdim:sec:optimiser}).
% src: research_benchmark/results/tables.md (T3: sobol/random after Adam -> L-BFGS 0.97, 0.54, 1.00, 0.99, 0.78; 10 of 10 seeds on no problem)
That the point set matters mainly when points are scarce is in the direction reported by
\citet{wu2023}; at our much smaller budgets we do not reproduce their ranking of grid and
resampled points on every problem. That a PINN needs residual points next to the initial
data (C4) is consistent with the propagation hypothesis of \citet{daw2023} and with causal
training~\citep{wangsankaran2024}; the control experiment isolates it for one problem. The
late-time behaviour of the heat problem (C2) is an elementary consequence of a loss made of
absolute residuals; the time-slice analysis quantifies its size on this test problem. The cost gap to classical solvers agrees with \citet{grossmann2024}, and
the penalty bias is classical~\citep{babuska1973,muller2022}. The residual loss and the Ritz
loss have been compared on elliptic problems by \citet{chen2020comparison}, who found, as we
do, that neither dominates the other. The constant of (C7) is classical: it is
$\delta_1^{-1/2}$ for the first biharmonic Steklov eigenvalue $\delta_1$, Payne's bound for
convex domains~\citep{payne1968bounds}, as stated in~\citet[Lemma~5.3]{bucur2011first}, already
gives a constant independent of $d$, and the torsion function of the cube gives one that
decreases slowly with $d$; to our knowledge, what is new is the explicit two-sided bound of
the squared constant $1/\delta_1$ of exact order $d^{-1/2}$ on the cube (equivalently, the
constant itself is of exact order $d^{-1/4}$). The error bound that
follows is an application of known theory with a better constant: a bound of the same form,
on smooth domains and with a constant that is not explicit, is in
\citet[App.~B]{shin2023error}. The remedies of (C8) use established
constructions~\citep{lagaris1998,ainsworth2021galerkin,chang2022horderdnn,zhang2025khorderdnn};
the paired comparisons measure their benefits and failure cases under a common budget.
The empirical contribution is the combination of
exact solutions, seeds, paired comparisons and a classical baseline of the kind whose
absence \citet{mcgreivy2024} criticise; a verification of the benchmark and of the dimension
study by re-aggregation of the raw runs and by separately written code for part of them
(Section~\ref{s3_methods:sec:verif} states what it covers and what it does not); and the
release of all code and data.

The experiments cover linear, constant-coefficient problems on boxes; the
networks are small and use the tanh activation; the loss weights are equal to one in
Sections~\ref{s4_lowdim:sec} and~\ref{s5_hard:sec}, and a single boundary weight per problem
and method is chosen once in Section~\ref{s6_highdim:sec}. Of the published remedies only two
cheap ones were tried, at $d=20$; adaptive loss weights, causal training, Fourier features and
boundary conditions built into the trial function were not
(Section~\ref{s8_discussion:sec:nottried}). Each of these might change a conclusion, in
particular the outcome on \PWone:
% src: refs.bib (rathore2024, annote: App. A.3, beta = 5; Sec. 2.2: widths up to 400, 10000 residual points, tuned Adam rate, switch after 1000, 11000 or 31000 Adam iterations; Table 1: lowest loss over widths after tuning)
on a variant with a faster second mode, \citet{rathore2024} report a few per cent with a
larger network and budget and tuned settings, without any of these remedies
(Section~\ref{s5_hard:sec:w1}).

Section~\ref{s2_problems:sec} states the problems and Section~\ref{s3_methods:sec} the
methods. Sections~\ref{s4_lowdim:sec}, \ref{s5_hard:sec} and~\ref{s6_highdim:sec} report
the low-dimensional benchmark, the harder problems with the finite-difference baselines,
and the dimension study, including the stability constant with the computable error bound
(Section~\ref{s6_highdim:sec:bound}) and the two remedies
(Section~\ref{s6_highdim:sec:remedies}). Section~\ref{s7_pitfalls:sec} demonstrates the two
implementation pitfalls, and Section~\ref{s8_discussion:sec} discusses the evidence and its
limitations. The Supplementary Material contains the proofs
(Section~\ref{sA_proofs:sec}), further results, figures and tables for each section, and
implementation details with full tables (Section~\ref{sC_details:sec}); its sections,
equations, figures and tables carry the prefix S.
